%! TeX program = lualatex
\documentclass[9pt]{beamer}

\usepackage{enumitem}                                  % list config
\usepackage{amssymb,amsmath,amsthm,stmaryrd,mathtools} % fonts, symbols
\usepackage{unicode-math}                              % fonts
\usepackage{tikz-cd}                                   % diagrams
\usepackage{hyperref}
\usepackage{booktabs}                                  % cmidrules
\usepackage{newunicodechar}
\usepackage[backend=biber,style=alphabetic]{biblatex}  % citations

\usepackage{adjustbox}

\usetheme{Montpellier}
\setbeamertemplate{page number in head/foot}[totalframenumber]
\setbeamertemplate{footline}[page number]

\setmainfont{fbb}
\setmathfont{Libertinus Math}
%FIXME: Trouver mieux que STIX Two Math
\setmathfont{STIX Two Math}[range={\obslash, \oslash, \sharp, \flat, \natural}]

% Shortcuts for math fonts
\newcommand{\bb}[1]{\mathbb{#1}}
\newcommand{\ca}[1]{\mathcal{#1}}
\newcommand{\fk}[1]{\mathfrak{#1}}
\newcommand{\sr}[1]{\mathscr{#1}}
\newcommand{\cat}[1]{\textcat{#1}}
\newcommand{\textcat}[1]{{\normalfont\textsf{#1}}}

% Utilities
\newcommand{\flip}[1]{\text{\rotatebox[origin=c]{-180}{$#1$}}}

% Shortcuts
% FIXME: blacktriangleleft and right too short, cf setmathfont
\newcommand{\PP}{{\bb P}}
\newcommand{\To}{\Rightarrow}
\newcommand{\coce}[2]{{\hat{#1}}_{#2}}
\newcommand{\coc}[2]{{\widehat{#1}}(#2)}
\newcommand{\colimeq}{\mathrel{\hat\approx}}
\newcommand{\colimleq}{\mathrel{\hat\lesssim}}
\newcommand{\copsh}[2]{{\widebreve{#1}}[{#2}]}
\newcommand{\dfrom}{\nvleftarrow}
\newcommand{\from}{\leftarrow}
\newcommand{\ip}{\flip\pi}
\newcommand{\iset}[1]{\llbracket #1 \rrbracket}
\newcommand{\lext}{\mathbin{\rhd}}
\newcommand{\lift}{\mathbin{\lhd}}
\newcommand{\nvequal}{=\!\!\!|\!\!\!=}
\newcommand{\op}{\mathrm{op}}
\newcommand{\oy}{\flip\yo}
\newcommand{\plext}{\mathbin{{\rhd}\mathclap{\mspace{-17.5mu}\cdot}}}
\newcommand{\plift}{\mathbin{\flip{{\rhd}\mathclap{\mspace{-17.5mu}\cdot}}}}
\newcommand{\prext}{\mathbin{{\blacktriangleright}\mathclap{\mspace{-10.5mu}\raisebox{-0.45pt}{\textcolor{white}{$\cdot$}}}}}
\newcommand{\prift}{\mathbin{\flip{{\blacktriangleright}\mathclap{\mspace{-10.5mu}\raisebox{-0.45pt}{\textcolor{white}{$\cdot$}}}}}}
\newcommand{\psh}[2]{{\widehat{#1}}[{#2}]}
\newcommand{\rext}{\mathbin{\blacktriangleright}}
\newcommand{\rift}{\mathbin{\blacktriangleleft}}
\newcommand{\set}[1]{\left\{#1\right\}}
\newcommand{\simeqq}{\cong}
\newcommand{\tp}[1]{\langle#1\rangle}
\newcommand{\unit}{{\tp{}}}

% Operators
\DeclareMathOperator{\Alg}{\cat{Alg}}
\DeclareMathOperator{\Cart}{\cat{Cart}}
\DeclareMathOperator{\Dist}{\bb{D}\textbf{ist}}
\DeclareMathOperator{\Env}{\cat{Env}}
\DeclareMathOperator{\Fib}{Fib}
\DeclareMathOperator{\Fun}{Fun}
\DeclareMathOperator{\Hom}{Hom}
\DeclareMathOperator{\Ind}{Ind}
\DeclareMathOperator{\Lan}{Lan}
\DeclareMathOperator{\Mod}{\bb{M}\textbf{od}}
\DeclareMathOperator{\Opalg}{\cat{Opalg}}
\DeclareMathOperator{\Pro}{Pro}
\DeclareMathOperator{\Sh}{\cat{Sh}}
\DeclareMathOperator{\Span}{\bb{S}\textbf{pan}}
\DeclareMathOperator{\Spec}{Spec}
\DeclareMathOperator{\bSpan}{\ca{S}\mathit{pan}}
\DeclareMathOperator{\cart}{cart}
\DeclareMathOperator{\coCart}{\cat coCart}
\DeclareMathOperator{\cod}{cod}
\DeclareMathOperator{\colim}{colim}
\DeclareMathOperator{\co}{co}
\DeclareMathOperator{\dom}{dom}
\DeclareMathOperator{\hFun}{hFun}
\DeclareMathOperator{\id}{id}
\DeclareMathOperator{\ob}{ob}
\DeclareMathOperator{\opcart}{opcart}

\DeclarePairedDelimiter\ceil{\lceil}{\rceil}
\DeclarePairedDelimiter\floor{\lfloor}{\rfloor}

% Adjunctions, horizontal arrows, pullback
\usetikzlibrary{decorations.pathmorphing}
\usetikzlibrary{decorations.markings}
\tikzset{%
    adj/.style={%
        /tikz/commutative diagrams/phantom,
        "\dashv"{rotate=-90}
    },
    loose/.default=0.5ex,
    loose/.style={%
        /utils/exec=\tikzset{every node/.append style={outer sep=0.8ex}},
        postaction=decorate,
        decoration={markings, mark=at position 0.5 with {\draw[-] (0,#1) -- (0,-#1);}}
    },
    pullback/.style={
        /tikz/commutative diagrams/phantom,
        "\scalebox{1.5}{$\lrcorner$}"{pos=0.125}
    }
}

% `よ`
\DeclareFontFamily{U}{min}{}
\DeclareFontShape{U}{min}{m}{n}{<-> udmj30}{}
\newcommand\yo{\!\text{\usefont{U}{min}{m}{n}\symbol{'207}}\!}

% widearc and widebreve
\makeatletter
\def\widebreve{\mathpalette\wide@breve}
\def\wide@breve#1#2{\sbox\z@{$#1#2$}%
     \mathop{\vbox{\m@th\ialign{##\crcr
\kern0.08em\brevefill#1{0.8\wd\z@}\crcr\noalign{\nointerlineskip}%
                    $\hss#1#2\hss$\crcr}}}\limits}
\def\brevefill#1#2{$\m@th\sbox\tw@{$#1($}%
  \hss\resizebox{#2}{\wd\tw@}{\rotatebox[origin=c]{90}{\upshape(}}\hss$}
\makeatother
\makeatletter
\def\widearc{\mathpalette\wide@arc}
\def\wide@arc#1#2{\sbox\z@{$#1#2$}%
     \mathop{\vbox{\m@th\ialign{##\crcr
\kern0.08em\arcfill#1{0.8\wd\z@}\crcr\noalign{\nointerlineskip}%
                    $\hss#1#2\hss$\crcr}}}\limits}
\def\arcfill#1#2{$\m@th\sbox\tw@{$#1($}%
  \hss\resizebox{#2}{\wd\tw@}{\rotatebox[origin=c]{-90}{\upshape(}}\hss$}
\makeatother

% Bijection of cells
\newenvironment{iffseq}{%
    \global\let\externaldblbackslash\\
    \begin{array}{cl}
    \ifundef{\internaldblbackslash}{%
        \global\let\internaldblbackslash\\%
        \gdef\\{\internaldblbackslash\cmidrule{1-1}\morecmidrules\cmidrule{1-1}}%
    }{}
}{%
    \end{array}
    \global\undef\internaldblbackslash
    \global\let\\\externaldblbackslash
}

\definecolor{ocre}{RGB}{0, 113, 188}
\hypersetup{
    colorlinks,
    citecolor=ocre,
    filecolor=ocre,
    linkcolor=ocre,
    urlcolor=ocre
}

% lists (fix babel)
\setlist[itemize]{label=\textbullet}
\setlist[enumerate]{label=(\alph*)}

% citations
\addbibresource{slides.bib}

\title{Théorie des catégories formelle}
\author{Luc Saccoccio--Le Guennec\\[1ex] {\small Encadrant : Morgan Rogers}}
\date{Mardi 22 Septembre 2026}

\begin{document}

\frame{\titlepage}

\begin{frame}{Théorie des catégories formelle}
    \begin{columns}[T]
        \begin{column}{0.5\textwidth}
        Plusieurs saveurs :
        \begin{itemize}
            \item Catégories enrichies (dans un cosmos, une bicatégorie, une catégorie virtuelle double...)
            \item Catégories (localement) internes
            \item Catégories fibrées et indexées
        \end{itemize}
        \end{column}
        \vrule{}
        \begin{column}{0.5\textwidth}
        Mais résultats communs
        \begin{itemize}
            \item Adjonctions
            \item Monades
            \item Yoneda
        \end{itemize}
        \end{column}
    \end{columns}
    \vspace{\baselineskip}
    \hrule
    \vspace{\baselineskip}
    \begin{itemize}
        \item Structures de Yoneda sur 2-catégories
        \item 2-monades lax-idempotentes (KZ-doctrines)
        \item Équipements de proflèches
    \end{itemize}
    Imparfaites: catégories enrichies sur une base quelconque, constructions des préfaisceaux, composition horizontale...
\end{frame}

\begin{frame}{Catégories $\cat V$-enrichies}
    On fixe une catégorie monoïdale symétrique close complète et cocomplète $\cat V$. Tout au long de l'exposé on instanciera les définitions et résultats dans les catégories $\cat V$-enrichies.
    \vspace{\baselineskip}
    \[\cat{Set}, \cat{sSet}, \cat{Cat}...\]
\end{frame}

\begin{frame}{Catégories virtuelles doubles}
    \begin{columns}[T]
        \begin{column}{0.5\textwidth}
            Une \textbf{catégorie virtuelle double} $\bb X$ est la donnée de :
            \begin{enumerate}
                \item Objets $A, B, C, \cdots$
                \item Flèches composables $\begin{tikzcd} X \ar[d, "f"]\\ Y \end{tikzcd}$
                \item Distributeurs $p : X \dfrom Y$
                \item Cellules
                    \[\adjustbox{scale=0.8,center}{\begin{tikzcd}[row sep=tiny,ampersand replacement=\&]
                X_0 \ar[dd, "f"'] \& \ar[l, "p_1"', loose] X_1
                                     \& \ar[l, "p_2"', loose] \cdots
                                     \& \ar[l, "p_{n-1}"', loose] X_{n-1}
                                     \& \ar[l, "p_n"', loose] X_n \ar[dd, "g"]\\
                \&\& \phi \&\& \\
                Y \&\&\&\& \ar[llll, "q", loose] Z
            \end{tikzcd}}\]
            \end{enumerate}
        \end{column}
        \vrule{}
        \begin{column}{0.5\textwidth}
            $\bb X = \cat V-\Dist$
            \begin{enumerate}
                \item $\cat V$-catégories
                \item $\cat V$-foncteurs
                \item $\cat V$-distributeurs
                    \[\adjustbox{scale=0.8}{\begin{tikzcd}
                            p : X^\op \otimes Y \to \cat V
                        \end{tikzcd}
                    }\]
                \item $\cat V$-transformations naturelles
                    \[\adjustbox{scale=0.8,center}{
                            \begin{tikzcd}[ampersand replacement=\&]
                                p_1(x_0, x_1) \otimes \cdots \otimes p_n(x_{n-1}, x_n) \ar[r, "{\phi_{x_0, \cdots, x_n}}"] \& q(fx_0, gx_n)
                            \end{tikzcd}
                    }\]
            \end{enumerate}
        \end{column}
    \end{columns}
\end{frame}

\begin{frame}{Cellules globulaires}
    On notera $\phi : p_1, \cdots, p_n \To q$ une cellule de la forme
    \[\adjustbox{scale=0.8}{\begin{tikzcd}[ampersand replacement=\&]
                X_0 \ar[dd, equal] \& \ar[l, "p_1"', loose] X_1
                                     \& \ar[l, "p_2"', loose] \cdots
                                     \& \ar[l, "p_{n-1}"', loose] X_{n-1}
                                     \& \ar[l, "p_n"', loose] X_n \ar[dd, equal]\\
                \&\& \phi \&\& \\
                Y \&\&\&\& \ar[llll, "q", loose] Z
        \end{tikzcd}}\]
    On notera aussi $\vec p$ pour $p_1, \cdots, p_n$.
\end{frame}

\begin{frame}{Cellules (op)cartésiennes}
    \begin{columns}[T]
        \begin{column}{0.5\textwidth}
            \begin{itemize}
                \item Cellule opcartésienne : $\phi : \vec q \To q$ avec une bijection
                    \[\begin{iffseq}
                        \vec p, \vec q, \vec r \To s\\
                        \vec p, q, \vec r \To s
                    \end{iffseq}\]
                    On note $q = q_1 \odot \cdots \odot q_n$.
                \item Cellule cartésienne :
                    $\adjustbox{scale=0.8}{\begin{tikzcd}[ampersand replacement=\&]
                        \ar[d, "f"'] A \& \ar[l, loose, "{q(f, g)}"', ""{anchor=center,name=t}] B \ar[d, "g"]\\
                        X \& \ar[l, loose, "q", ""{anchor=center,name=b}] Y
                        \ar[from=t, to=b, phantom, "\cart"{anchor=center}]
                    \end{tikzcd}}$
                    \[\hspace{-2em}\adjustbox{scale=0.8}{\begin{tikzcd}[row sep=tiny, ampersand replacement=\&]
                        \cdot \ar[dd, "gf"'] \& \ar[l, "r_1"', loose] \cdots \& \ar[l, "r_n"', loose] \cdot \ar[dd, "ih"]\\
                                            \& \phi \& \\
                        \cdot \&\& \ar[ll, "q", loose] \cdot
                    \end{tikzcd}}
                    \hspace{0.5em}=\hspace{0.5em}
                    \adjustbox{scale=0.8}{\begin{tikzcd}[row sep=tiny, ampersand replacement=\&]
                        \cdot \ar[dd, "f"'] \& \ar[l, "r_1"', loose] \cdots \& \ar[l, "r_n"', loose] \cdot \ar[dd, "h"]\\
                                           \& \tilde{\phi} \& \\
                        \cdot \ar[dd, "g"'] \&\& \ar[ll, "p", loose] \cdot \ar[dd, "i"]\\
                                            \& \text{cart} \&\\
                        \cdot \&\& \ar[ll, "q", loose] \cdot
                    \end{tikzcd}}\]
            \end{itemize}
        \end{column}
        \vrule{}
        \begin{column}{0.5\textwidth}
            Dans $\bb X = \cat V-\Dist$
            \begin{itemize}
                \item $p : B \dfrom C, q : C \dfrom D$
                    \[(p\odot q)(b, d) \simeqq \int^{c\in\ob C} p(b, c)\otimes q(c, d)\]
                \item $q(f, g)(a, b) = q(fa, gb)$
            \end{itemize}
        \end{column}
    \end{columns}
\end{frame}

\begin{frame}{Équipements virtuels}
    \begin{columns}[T]
        \begin{column}{0.5\textwidth}
            Un équipement virtuel a les
            \begin{itemize}
                \item Distributeurs identités $A(1, 1)$
                    \[\adjustbox{scale=0.8}{\begin{tikzcd}[ampersand replacement=\&]
                        \& \ar[dl, equal] A \ar[dr, equal] \&\\
                        A \&{}\& \ar[ll, loose, equal, "{A(1, 1)}"] A
                        \ar[from=1-2, to=2-2, phantom, "\opcart"{anchor=center}]
                    \end{tikzcd}}\]
                \item Restrictions $q(f, g)$
                    \[\adjustbox{scale=0.8}{\begin{tikzcd}[ampersand replacement=\&]
                            \ar[d, "f"'] A \& \ar[l, loose, "{q(f, g)}"', ""{anchor=center,name=t}] B \ar[d, "g"]\\
                            X \& \ar[l, loose, "q", ""{anchor=center,name=b}] Y
                            \ar[from=t, to=b, phantom, "\cart"{anchor=center}]
                    \end{tikzcd}}\]
            \end{itemize}
        \end{column}
        \vrule{}
        \begin{column}{0.5\textwidth}
            Dans $\bb X = \cat V-\Dist$
            \begin{itemize}
                \item \(\begin{aligned}
                        A(1, 1)(x, y) &= A(x, y)\\
                                     &:=\Hom_A(x, y)
                    \end{aligned}\)
                \item $q(f, g)(a, b) = q(fa, gb)$
            \end{itemize}
            \vspace{\baselineskip}
            \hrule
            \vspace{\baselineskip}
            \[A(f, g) := A(1, 1)(f, g)\]
            On fera un \textit{choix} de restrictions fonctorielles. Un tel équipement virtuel est \textbf{strict}.
            \begin{align*}
                p(1, 1) &= p & p(f, g)(h, i) &= p(fh, gi)
            \end{align*}
        \end{column}
    \end{columns}
\end{frame}

\begin{frame}{Relèvements et extensions}
    \begin{columns}[T]
        \begin{column}{0.5\textwidth}
            \begin{enumerate}
                \item Relèvement à droite
                    $\adjustbox{scale=0.8}{\begin{tikzcd}[ampersand replacement=\&]
                            Y \ar[dr, "\vec{p}"', ""{name=p, anchor=center}, loose] \& \ar[l, loose, "q \rift \vec{p}"'] X \ar[d, loose, "q", ""{name=q, anchor=center}]\\
                            \ar[from=p, to=q, Rightarrow, shift left=2, xshift=-1mm]
                                                      \& Z
                        \end{tikzcd}}$ avec
                        \begin{iffseq}
                            \vec r \To q\rift\vec p\\
                            \vec p, \vec r \To q
                        \end{iffseq}
                \item Extension à droite
                    $\adjustbox{scale=0.8}{\begin{tikzcd}[ampersand replacement=\&]
                            Y \& \ar[l, loose, "\vec p \rext q"'] Z\\
                            \& \ar[ul, loose, "q", ""{anchor=center,name=q}] X \ar[u, loose, "\vec p"']
                            \ar[from=1-2, to=q, Rightarrow]
                        \end{tikzcd}}$ avec
                        \begin{iffseq}
                            \vec r \To \vec p \rext q\\
                            \vec r, \vec p \To q
                        \end{iffseq}
            \end{enumerate}
        \end{column}
        \vrule{}
        \begin{column}{0.5\textwidth}
            Dans $\bb X = \cat V-\Dist$
            \begin{align*}
                (q\rift p)(x, y) &\simeqq \ca PA(p(-, x), q(-, y))\\
                                 &\simeqq \int_{z\in\ob Z}\cat V(p(z, x) q(z, y))
            \end{align*}

            \vspace{\baselineskip}
            En particulier, n'existe pas nécessairement si $\cat V$ n'est pas fermé.
        \end{column}
    \end{columns}
\end{frame}

\begin{frame}{Colimites et extensions de Kan}
    \begin{columns}[T]
        \begin{column}{0.5\textwidth}
            \begin{enumerate}
                \item Colimites pondérées : $Y(\colim^{\vec p}f, 1) \simeq Y(f, 1) \rift \vec p$
                    \[\adjustbox{scale=0.8}{\begin{tikzcd}[ampersand replacement=\&]
                            X \ar[dr, loose, "\vec p"'] \ar[r, "\colim^{\vec p}f"] \& Y\\
                            \& A \ar[u, "f"']
                    \end{tikzcd}}\]
                \item Extension de Kan à gauche : $j\plext f \simeqq \colim^{E(j, 1)}f\simeqq Y(f, 1) \rift E(j, 1)$
                    \[\adjustbox{scale=0.8}{\begin{tikzcd}[ampersand replacement=\&]
                            E \ar[r, "j\plext f"] \& Y\\
                            \& \ar[ul, "j"] A \ar[u, "f"']
                    \end{tikzcd}}\]
            \end{enumerate}
        \end{column}
        \vrule{}
        \begin{column}{0.5\textwidth}
            Dans $\bb X = \cat V-\Dist$
            \begin{itemize}
                \item $Y((\colim^pf)x, y) \simeqq \ca PA(p(-, x), Y(f-, y))$
                \item $Y((j\plext f)x, y) \simeq \ca PA(E(j-, x), Y(f-, y))$
            \end{itemize}
            \vspace{\baselineskip}
            \hrule
            \vspace{\baselineskip}
            On peut obtenir la caractérisation 2-catégorique :
            \begin{iffseq}
                j\plext f \To x\\
                f \To x\circ j
            \end{iffseq}
        \end{column}
    \end{columns}
\end{frame}

\begin{frame}{Que peut-on étudier ?}
    \begin{enumerate}
        \item Densité et pleine fidélité de flèches.
        \item (co)Limites absolues.
        \item Monades et adjonctions relatives.
        \item Algèbres et algèbres libres.
        \item Cocomplétions.
        \item Objets en préfaisceaux.
        \item Dualité d'Isbell et point fixe.
    \end{enumerate}
    \vspace{\baselineskip}
    \hrule
    \vspace{\baselineskip}
    \begin{itemize}
        \item Préservation, réflexion, relèvement et création de (co)limites.
        \item Monadicité relative.
    \end{itemize}
\end{frame}

\begin{frame}{Densité et pleine fidélité}
    \begin{columns}[T]
        \begin{column}{0.5\textwidth}
            \begin{enumerate}
                \item Pleine fidélité (f.f.) : $B(f, f) \simeqq A(1, 1)$
                \item Densité : $p\rift p \simeqq 1_B$ et $f\plext f \simeqq \id_B$.
                    \[\adjustbox{scale=0.8}{\begin{tikzcd}[ampersand replacement=\&]
                            A \& \ar[l, loose, "p"'] B \ar[d, equal, loose]\\
                            \& \ar[ul, loose, "p"] B
                    \end{tikzcd}}
                    \adjustbox{scale=0.8}{\begin{tikzcd}[ampersand replacement=\&]
                            A \ar[r, "f"] \ar[dr, "f"] \& B \ar[d, equal]\\
                            \& B
                    \end{tikzcd}}\]
            \end{enumerate}
        \end{column}
        \vrule{}
        \begin{column}{0.5\textwidth}
        Dans $\bb X = \cat V-\Dist$
            \begin{enumerate}
                \item $\forall c, d, B(fc, fd) \simeqq A(c, d)$.
                \item $\Lan_ff$ existe et est isomorphe à $\id_B$.
            \end{enumerate}
        \end{column}
    \end{columns}
    \vspace{\baselineskip}
    \hrule
    \vspace{\baselineskip}
    La densité se définit horizontalement et se récupère grâce au conjoint. On a la caractérisation équivalente
    \begin{columns}[T]
        \begin{column}{0.5\textwidth}
            \[\begin{iffseq}
                \vec r \To E(g, h)\\
                p(1, g), \vec r \To p(1, h)
            \end{iffseq}\]
        \end{column}
        \begin{column}{0.5\textwidth}
            \[\begin{iffseq}
                g \To h\\
                p(1, g) \To p(1, h)
            \end{iffseq}\]
        \end{column}
    \end{columns}
\end{frame}

\begin{frame}{$j$-Absoluité}
    Pour $j : A \to E$, $\colim^{\vec p}f$ est $j$-absolue si
    $$E(j, \colim^{\vec p}f) \simeqq E(j, f) \odot_L \vec p$$
    $\colim^{\vec p}f$ est préservé par tout les $j\plext g$, en particulier tout les foncteurs préservent les colimites $1_E$-absolues.
    \vspace{\baselineskip}
    \pause
    \hrule
    \vspace{\baselineskip}
    \begin{itemize}
        \item Dans $\cat{Set}-\Dist$, les coégaliseurs split ou les pushouts dans une catégorie de Reedy élégante sont absolus.
        \item Dans $\cat{Ab}-\Dist$, les biproduits initiaux sont absolus.
        \item Dans $\overline{\bb R_+}-\Dist$, les limites de suites de Cauchy sont absolues.
    \end{itemize}
\end{frame}

\begin{frame}{Adjonctions et monades horizontales}
    Postulant l'existence des compositions, on peut définir les notions horizontalement :
    \begin{columns}[T]
        \begin{column}{0.5\textwidth}
            \[\ell \dashv r\]
            \begin{enumerate}
                \item $\ell : C \dfrom A$, $r : A \dfrom C$.
                \item $\eta : \To r\odot\ell$.
                \item $\varepsilon : \ell, r\To C(1, 1)$.
            \end{enumerate}
            \begin{align*}
                (1_\ell, \varepsilon)(e_\ell, \eta) &= e_\ell & (\varepsilon, 1_r)(\eta, e_r) = e_r
            \end{align*}
        \end{column}
        \vrule{}
        \begin{column}{0.5\textwidth}
            \[(t, \mu, \eta)\]
            \begin{enumerate}
                \item $t : A \dfrom A$
                \item $\mu : t, t\To t$
                \item $\eta : \To t$
            \end{enumerate}
            \begin{align*}
                \mu(\mu, e_t) &= \mu(e_t, \mu) & \mu(\eta, e_t) &= \mu(e_t, \eta) = e_t
            \end{align*}
        \end{column}
    \end{columns}
    \pause
    \vspace{\baselineskip}
    \hrule
    \vspace{\baselineskip}
    Lien avec
    \begin{itemize}
        \item les flèches pleinement fidèles,
        \item les adjonctions relatives,
        \item l'équipement virtuel des monades $\Mod(\bb X)$
        \item ...
    \end{itemize}
\end{frame}

\begin{frame}{Adjonctions relatives}
    \begin{columns}[T]
        \begin{column}{0.5\textwidth}
            \begin{enumerate}
                \item $j : A \to E$
                \item $\ell : A \to C$
                \item $r : C \to E$
                \item $\sharp : C(\ell, 1) \simeqq E(j, r) : \flat$
            \end{enumerate}
            \[\adjustbox{scale=0.8}{\begin{tikzcd}[ampersand replacement=\&]
                \& C \ar[dr, "r", ""{anchor=north,name=radj}] \&\\
                A \ar[ur, "\ell", ""{anchor=north,name=ladj}] \ar[rr, "j"']  \&\& E
                \ar[from=ladj, to=radj, phantom, "\dashv"{anchor=center}]
        \end{tikzcd}}\]
        De façon équivalente :
        \begin{enumerate}
            \item $\eta : j \To r\circ \ell$
            \item $\varepsilon : C(1, \ell), E(j, r) \To C(1, 1)$
        \end{enumerate}
        \end{column}
        \vrule{}
        \begin{column}{0.5\textwidth}
            On retrouve$^*$ :
            \begin{enumerate}
                \item Essentielle unicité des adjoints.
                \item Préservation des (co)limites.
                \item Construction des adjoints comme extension ou relèvement de Kan.
                \[\adjustbox{scale=0.8}{\begin{tikzcd}[ampersand replacement=\&]
                    \&  C \ar[dr, "r"] \&\\
                    A \ar[ur, dashed, "j\plift r"] \ar[rr, "j"'] \&{}\& E
                    \ar[from=1-2, to=2-2, phantom, "\dashv"{anchor=center}]
                \end{tikzcd}}\]
            \end{enumerate}
            \vspace{\baselineskip}
            \hrule
            \vspace{\baselineskip}
            \[C(1, \ell) \dashv E(j, r)\]
        \end{column}
    \end{columns}
\end{frame}

\begin{frame}{Constructions}
    \begin{columns}[T]
        \begin{column}{0.5\textwidth}
            Précomposition :
            \[\adjustbox{scale=0.8}{\begin{tikzcd}[ampersand replacement=\&]
                    \&\&C\ar[dr, "r"]\&\\
                    A \ar[r, "\ell'"] \& B \ar[ur, "\ell"] \ar[rr, "j"'] \&{}\& D
                    \ar[from=1-3, to=2-3, phantom, "\dashv"{anchor=center}]
                \end{tikzcd}}\]
                \[\adjustbox{scale=0.8}{\begin{tikzcd}[ampersand replacement=\&]
                    \&C\ar[dr, "r"]\&\\
                    A \ar[ur, "\ell\ell'"] \ar[rr, "j\ell'"'] \&{}\& D
                    \ar[from=1-2, to=2-2, phantom, "\dashv"{anchor=center}]
                \end{tikzcd}}\]
        \end{column}
        \vrule{}
        \begin{column}{0.5\textwidth}
            Collage :
            \[\adjustbox{scale=0.8}{\begin{tikzcd}[ampersand replacement=\&]
                \&C \ar[dr, "r"] \&\&\\
                \&\& \ar[dr, "r'", ""{anchor=center,name=r_}] D \&\\
                A \ar[uur, "\ell"] \ar[urr, "\ell'", ""{anchor=center,name=l_}] \ar[rrr, "j"'] \&\&\& E\\
                \ar[from=l_, to=r_, phantom, "\dashv", shift right]
            \end{tikzcd}}\]
        \end{column}
    \end{columns}
    \vspace{\baselineskip}
    \hrule
    \vspace{\baselineskip}
    Postcomposition :
    \[\adjustbox{scale=0.8}{\begin{tikzcd}[ampersand replacement=\&]
        \&C\ar[dr, "r"]\&\&\\
        A \ar[ur, "\ell"] \ar[rr, "j"'] \&{}\& B \ar[r, "r'", hook] \& D
        \ar[from=1-2, to=2-2, phantom, "\dashv"{anchor=center}]
    \end{tikzcd}}
    \hspace{0.5em}
    \implies
    \hspace{0.5em}
    \adjustbox{scale=0.8}{\begin{tikzcd}[ampersand replacement=\&]
        \&C\ar[dr, "r'r"]\&\\
        A \ar[ur, "\ell"] \ar[rr, "r'j"'] \&{}\& D
        \ar[from=1-2, to=2-2, phantom, "\dashv"{anchor=center}]
    \end{tikzcd}}\]
\end{frame}

\begin{frame}{Monades relatives}
    \begin{columns}[T]
        \begin{column}{0.5\textwidth}
            \[T := (t, \dag, \eta)\]
            \begin{enumerate}
                \item $j : A \to E$
                \item $t : A \to E$
                \item $\dag : E(j, t) \To E(t, t)$
                \item $\eta : j \To t$
            \end{enumerate}
            Source : Adjonctions relatives
        \end{column}
        \vrule{}
        \begin{column}{0.5\textwidth}
            \begin{center}
                Constructions
            \end{center}
            \begin{enumerate}
                \item $C(1, \ell) \dashv E(j, r)$ induit $E(j, T)$.
                \item Précomposition
                \item Collage
                    \[\adjustbox{scale=0.8}{\begin{tikzcd}[ampersand replacement=\&]
                            \& D \ar[dr, "r"] \&\\
                            A \ar[ur, "\ell"{description}, ""{anchor=center,name=ell}] \ar[ur, "t", bend left=50, ""{anchor=center,name=t}] \ar[rr, "j"'] \&{}\& E
                            \ar[from=ell, to=t, Rightarrow, "\eta"'] \ar[from=1-2, to=2-2, phantom, "\dashv"{anchor=center}]
                    \end{tikzcd}}\]
            \end{enumerate}
        \end{column}
    \end{columns}
    \vspace{\baselineskip}
    \hrule
    \vspace{\baselineskip}
    \[\adjustbox{scale=0.8}{\begin{tikzcd}[ampersand replacement=\&]
            \& B \ar[r, "j"{description}] \ar[r, "t", bend left=50] \& D \ar[dr, "r", ""{anchor=center,name=r_}] \&\\
            A \ar[ru, "\ell"] \ar[rru, "j\ell"{description}, ""{anchor=center,name=jl}] \ar[rrr, "j"'] \&\&\& E
            \ar[from=jl, to=r_, phantom, "\dashv"{anchor=center}]
    \end{tikzcd}}\]
\end{frame}

\begin{frame}{Algèbres}
    \[T = (t, \dag, \eta)\]
    \begin{columns}[T]
        \begin{column}{0.5\textwidth}
            \begin{enumerate}
                \item $D$
                \item $e : D \to E$
                \item $\rtimes : E(j, e) \To E(t, e)$
            \end{enumerate}
            \[E(\eta, e)\rtimes = e_{E(j, e)}\]
            \[\smallsmile_t(t, e)(\dag, \rtimes) = \rtimes\smallsmile_t(j, e)(e_{E(j, t)}, \rtimes)\]
            \vspace{\baselineskip}
            \hrule
            \vspace{\baselineskip}
            Morphisme $\vec p$-gradué $(e, \rtimes) \to (e', \rtimes')$
            \[\varepsilon : \vec p \To E(e, e')\]
        \end{column}
        \vrule{}
        \begin{column}{0.5\textwidth}
            \[(u_T : \Alg(T) \to E, \rtimes_T)\]
            \begin{enumerate}
                \item $(e, \rtimes)$ se factorise en $e = u_T\unit_{(e, \rtimes)}$ et $\rtimes = \rtimes_T(1, \unit_{(e, \rtimes)})$.
                \item $\varepsilon : E(1, e), \vec p \To E(1, e')$ se factorise en $\varepsilon = E(1, u_T), \unit_\varepsilon$ avec
            \end{enumerate}
            $$\unit_\varepsilon :\Alg(T)(1, \unit_{(e, \rtimes)}), \vec p \To \Alg(1, \unit_{(e', \rtimes')})$$
            \vspace{\baselineskip}
            \hrule
            \vspace{\baselineskip}
            \[\adjustbox{scale=0.8}{\begin{tikzcd}[ampersand replacement=\&]
                \& \Alg(T) \ar[dr, "u_T"] \&\\
                A \ar[ur, "f_T"] \ar[rr, "j"'] \&{}\& E
                \ar[from=1-2, to=2-2, phantom, "\dashv"{anchor=center}]
            \end{tikzcd}}\]
        \end{column}
    \end{columns}
\end{frame}

\begin{frame}{Résolutions}
    Soit $T$ admettant un objet en algèbre et un objet en opalgèbre.
    \begin{itemize}
        \item $f_T\dashv_ju_T$ est terminale.
        \item $k_T\dashv_jv_T$ est initiale.
        \item Il existe un unique $i_T : \Opalg(T) \to \Alg(T)$ tel que
        \[\begin{tikzcd}[ampersand replacement=\&]
            {\Opalg(T)} \ar[r, "i_T"] \& {\Alg(T)}\\
            A \ar[u, "k_T"] \ar[ur, "f_T"']
        \end{tikzcd}
        \hspace{3em}
        \begin{tikzcd}[ampersand replacement=\&]
            {\Opalg(T)} \ar[r, "i_T"] \ar[dr, "v_T"'] \& {\Alg(T)} \ar[d, "u_T"]\\
                                                      \& E
        \end{tikzcd}\]
    \end{itemize}
    \vspace{\baselineskip}
    \hrule
    \vspace{\baselineskip}
    Dans $\bb X = \cat V-\Dist$:
    \begin{columns}[T]
        \begin{column}{0.65\textwidth}
            \begin{itemize}
                \item Tout monade relative admet un objet en opalgèbres.
                \item Si $\cat V$ est fermée et admet les égaliseurs, tout $j$-monade admet un objet en algèbres.
            \end{itemize}
        \end{column}
        \vrule{}
        \begin{column}{0.35\textwidth}
            \begin{itemize}
                \item $\cat V$-Catégorie de Kleisli.
                \item $\cat V$-Catégorie d'Eilenberg-Moore.
            \end{itemize}
        \end{column}
    \end{columns}
\end{frame}

\begin{frame}{Objets en préfaisceaux}
    \begin{columns}[T]
        \begin{column}{0.5\textwidth}
            $A$ un objet, $\PP$ une classe de distributeurs
            \begin{enumerate}
                \item $\pi_A^\PP : A \dfrom \psh\PP A$ dense
                \item $\forall p : A \dfrom X \in \PP, \exists! \widearc p : X \to \psh\PP A, p = \pi_A^\PP(1, \widearc p)$
                \item $\forall f : X \to \psh\PP A, \pi_A^\PP(1, f)\in\PP$
            \end{enumerate}
        \end{column}
        \vrule{}
        \begin{column}{0.5\textwidth}
            Dans $\bb X = \cat V-\Dist$, $\PP = \ca P$. $\psh\PP A$ a:
            \begin{enumerate}
                \item Objets: $\cat V$-préfaisceaux sur $A$
                \item $\psh\PP A(f, g) = \int_{a\in A}\cat V(f(a), g(a))$
                \item $\pi_A^\PP(a, f) = f(a)$
            \end{enumerate}
        \end{column}
    \end{columns}
    \vspace{\baselineskip}
    \hrule
    \vspace{\baselineskip}
    Premières propriétés
    \begin{columns}[T]
        \begin{column}{0.75\textwidth}
            \begin{enumerate}
                \item $\psh\PP A(\widearc{p}, \widearc{q}) \simeqq q \rift p$
                \item $\colim^{\vec p}\widearc q$ existe si et seulement si $(q\odot\vec p)\in\PP^\simeqq$.
                \item Si $A(1, 1)$, $\yo_A^\PP := \widearc{A(1, 1)}$ et $\psh\PP A(\yo_A^\PP, 1) \simeqq \pi_A^\PP$.
                \item $\yo_A^\PP$ dense, pleinement fidèle et crée les limites pondérées.
            \end{enumerate}
        \end{column}
        \begin{column}{0.25\textwidth}
            \hspace{-5em}
            \begin{tikzcd}[ampersand replacement=\&]
                    \psh\PP A \ar[rr, hook, "\widearc{\pi_A^\PP}"] \&\& \psh{\PP'}A\\
                                                                   \& \ar[lu, "\yo_A^\PP"] A \ar[ru, "\yo_A^{\PP'}"']
            \end{tikzcd}
        \end{column}
    \end{columns}
\end{frame}

\begin{frame}{Cocomplétions}
    $A$ un objet, $\Phi$ une classe de poids
    \begin{enumerate}
        \item $\coce\Phi A : A  \to \coc\Phi A$
        \item $\coc\Phi A$ est $\Phi$-cocomplète
        \item Toute flèche $f : A \to X$ avec $X$ cocomplet se factorise en $f \simeqq \tilde{f} \circ \coce\Phi A$ avec $\tilde f  = \coce\Phi A \plext f : \coc\Phi A \to X$.
        \item Tout distributeur $p$ s'écrit $p \simeqq \tilde p(\coce\Phi A, 1)$ avec $\tilde p = p\rift\coc\Phi A(\coce\Phi A, 1) : \coc\Phi A \dfrom X$.
    \end{enumerate}
    \pause
    \vspace{\baselineskip}
    \hrule
    \vspace{\baselineskip}
    \begin{itemize}
        \item $\coce\Phi A : A \to \coc\Phi A$ est pleinement fidèle et dense.
        \item Si les $\tilde p \simeqq p\rift\coc\Phi A(\coce\Phi A, 1)$ sont représentables par $\widearc p$, cette flèche est essentiellement unique avec cette propriété : $\pi_A(1, \widearc p) \simeqq p$ avec $\pi_A := \coc\Phi A(\coce\Phi A, 1)$.
    \end{itemize}
    Une cocomplétion est \textit{presque} un objet en préfaisceaux. La distinction décèle le problème de \textbf{taille}.
\end{frame}

\begin{frame}{Cocomplétions comme objets en préfaisceaux}
    \begin{columns}[T]
        \begin{column}{0.5\textwidth}
            $\Phi$, $\PP$ et $A$ un objet $\PP$-admissible tels que
            \begin{enumerate}
                \item $\forall\vec p\in\Phi, q\in\PP, q\odot_L\vec p\in\PP^\simeqq$.
                \item Tout les relèvements à droite par $\pi_A^\PP : A \dfrom \psh\PP A$ existent.
                \item $\Phi$ induit les $\pi_A^\PP$-colimites.
            \end{enumerate}
            Alors $\yo_A^\PP : A \to \psh\PP A$ est la $\Phi$-cocomplétion de $A$.
        \end{column}
        \vrule{}
        \begin{column}{0.5\textwidth}
            Dans $\bb X = \cat V-\Dist$, $A$ petit et tel que tout $p\in\Phi$ a un petit codomaine.\\
            $\coc\Phi A \hookrightarrow \psh\PP A$ est la plus petite sous-$\cat V$-catégorie
            \begin{itemize}
                \item Contenant les représentables
                \item Close par $\Phi$-colimites.
            \end{itemize}
        \end{column}
    \end{columns}
    \vspace{\baselineskip}
    \hrule
    \vspace{\baselineskip}
    \begin{itemize}
        \item Si $\PP$ clôse par composées à gauche, $\psh\PP A$ est $\PP$-complet.
        \item Si $\pi_A^\PP$ relève à droite, $\yo_A^\PP$ exhibe la $\PP$-cocomplétion de $A$.
    \end{itemize}
\end{frame}

\begin{frame}{Dualité d'Isbell}
    Si $\psh\PP A(1, \yo_A^\PP), \copsh\PP A(\oy_A^\PP, 1)\in\PP^\simeqq$:
    \[\adjustbox{scale=0.8}{\begin{tikzcd}[ampersand replacement=\&]
            \ca{O} : \psh\PP A \ar[r, bend left=15, ""{name=O, below}] \& \ar[l, bend left=15, ""{name=Spec, above}] \copsh\PP A : \Spec \ar[from=O, to=Spec, adj]
    \end{tikzcd}}\]
    Avec $\ca O = \yo_A^\PP \plext \oy_A^\PP$ et $\Spec = \oy_A^\PP \prext \yo_A^\PP$.
    \[\adjustbox{scale=0.8}{\begin{tikzcd}[ampersand replacement=\&]
            \ar[dd, hook, "\widearc{\pi_A^\PP}"'] \psh\PP A \ar[rr, bend left=15, "\ca{O}", ""{anchor=center,name=O}] \&\& \ar[ll, bend left=15, "\Spec", ""{anchor=center,name=Spec}] \copsh\PP A \ar[dd, hook, "\widebreve{\ip_A^\PP}"] \\
            \& A \ar[ul, "\yo_A^\PP"] \ar[ur, "\oy_A^\PP"'] \ar[dl, "\yo_A^{\PP'}"'] \ar[dr, "\oy_A^{\PP'}"]\&\\
            \psh{\PP'} A \ar[rr, bend left=15, "\ca{O}'", ""{anchor=center,name=O_}] \&\& \ar[ll, bend left=15, "\Spec'", ""{anchor=center,name=Spec_}] \copsh{\PP'} A \\
            \ar[from=O, to=Spec, adj] \ar[from=O_, to=Spec_, adj]
    \end{tikzcd}}\]
\end{frame}

\begin{frame}{Point fixe de l'adjonction}
    D'après \cite{averyleinster}, adapté pour un équipement virtuel.
    \begin{enumerate}
        \item Une flèche $x_A : X \to \psh\PP A$ est réflexive si $\eta : A \to \Spec\circ\ca O$ induit un isomorphisme $\eta x_A : x_A \to \Spec\circ\ca O\circ x_A$.
        \item Une complétion reflexive est une flèche f.f. et  réflexive $\ca r_A^\PP : \ca R^\PP(A) \to \psh\PP A$ à travers laquelle toutes les flèches réflexives se factorisent.
            \begin{itemize}
                \item $\yo_A^\PP$ est réflexive
            \end{itemize}
    \end{enumerate}
    \vspace{\baselineskip}
    \hrule
    \vspace{\baselineskip}
    Problème : l'étude d'Avery et Leinster se base sur les \textbf{petits} $\cat V$-préfaisceaux.
\end{frame}

\begin{frame}{Propriété universelle du point fixe (1)}
    Une flèche est $\PP$-petite si son nerf $B(f, 1)$ est dans $\PP^\simeqq$. Elle est adéquate si elle est dense, codense et pleinement fidèle, et $\PP$-adéquate si elle est de plus $\PP$-petite et copetite.
    \begin{itemize}
        \item Stables par composition.
        \item $\ca J_A^\PP$ induite par $\yo_A^\PP$ est $\PP$-adéquate.
        \item Pour $f$ f.f., dense et $\PP$-petite, $n_f^\PP$ réflexive si et seulement si $f$ $\PP$-adéquate.
    \end{itemize}
    Pour $f$ $\PP$-adéquate, $\ca r_f^\PP$ induite par $n_f^\PP$ est adéquate et unique telle que
    \[\adjustbox{scale=0.8}{\begin{tikzcd}[ampersand replacement=\&]
            B \ar[rr, "\ca{r}_f^\PP"] \&\& \ca{R}^\PP(A)\\
            \& A \ar[ul, "f"] \ar[ur, hook, "\ca{J}_A^\PP"'] \&
    \end{tikzcd}}\]
\end{frame}

\begin{frame}{Propriété universelle du point fixe (2)}
    \begin{itemize}
        \item Cette construction est fonctorielle.
        \item $\ca r_f^\PP$ est $\PP$-adéquate si et seulement si $\ca R^\PP(f)$ l'est si et seulement si $\ca R^\PP(f)$ est une équivalence.
        \item Si on peut l'itérer, $\ca R^\PP(A) \simeq \ca R^\PP (\ca R^\PP(A))$.
        \item En particulier $\ca J_A^\PP$ est une équivalence si $A$ est une complétion réflexive.
        \item $\ca J_A^\PP$ préserve et réfléchit les (co)limites pondérées, $\ca r_A^\PP$ préserve les limites et réfléchit les (co)limites pondérées.
        \item $\ca R^\PP(A)$ est stable par les colimites absolues que $\psh\PP A$ admet.
    \end{itemize}
\end{frame}

\begin{frame}{Bibliographie}
    \nocite{*}
    \printbibliography[heading=bibintoc]
\end{frame}

\begin{frame}{Opalgèbres}
    \[T = (t, \dag, \eta)\]
    \begin{columns}[T]
        \begin{column}{0.5\textwidth}
            \begin{enumerate}
                \item $B$
                \item $a : A \to B$
                \item $\ltimes : E(j, t) \To B(a, a)$
            \end{enumerate}
            \[\smallfrown_a = \ltimes E(j, \eta)\smallfrown_j\]
            \[\ltimes\smallsmile_t(j, t)(e_{E(j, t)}, \dag) = \smallsmile_a(a, a)(\ltimes, \ltimes)\]
            \vspace{\baselineskip}
            \hrule
            \vspace{\baselineskip}
            Morphisme $\vec p$-gradué $(a, \ltimes) \to (a', \ltimes')$
            \[\alpha : \vec p, B(1, a) \To B'(1, a')\]
        \end{column}
        \vrule{}
        \begin{column}{0.5\textwidth}
            \[(k_T : A \to \Opalg(T), \ltimes_T)\]
            \begin{enumerate}
                \item $(a, \ltimes)$ se factorise en $a = []_{(a, \ltimes)}k_T$ et $\ltimes = []_{(a, \ltimes)}\ltimes_T$.
                \item $\alpha : \vec p, B(1, a) \To B(1, a')$ se factorise en $\alpha = []_\alpha(1, k_T)$ avec
            \end{enumerate}
            \[[]_\alpha : \vec p, \Opalg(T)(1, []_{(a, \ltimes)}) \To B'(1, []_{(a', \ltimes')})\]
            \vspace{\baselineskip}
            \hrule
            \vspace{\baselineskip}
            \[\adjustbox{scale=0.8}{\begin{tikzcd}[ampersand replacement=\&]
                \& \Opalg(T) \ar[dr, "v_T"] \&\\
                A \ar[ur, "k_T"] \ar[rr, "j"'] \&{}\& E
                \ar[from=1-2, to=2-2, phantom, "\dashv"{anchor=center}]
            \end{tikzcd}}\]
        \end{column}
    \end{columns}
\end{frame}

\end{document}
